\documentclass{article}

\usepackage{corl_2026} % Use this for the initial submission.
% \usepackage[final]{corl_2026} % Uncomment for the camera-ready ``final'' version.
% \usepackage[preprint]{corl_2026} % Uncomment for pre-prints (e.g., arxiv); This is like ``final'', but will remove the CORL footnote.

\title{Formatting Instructions for CoRL 2026 Camera-Ready}

% The \author macro works with any number of authors. There are two
% commands used to separate the names and addresses of multiple
% authors: \And and \AND.
%
% Using \And between authors leaves it to LaTeX to determine where to
% break the lines. Using \AND forces a line break at that point. So,
% if LaTeX puts 3 of 4 authors names on the first line, and the last
% on the second line, try using \AND instead of \And before the third
% author name.

% NOTE: authors will be visible only in the camera-ready and preprint versions (i.e., when using the option 'final' or 'preprint'). 
% 	For the initial submission the authors will be anonymized.

\author{
  Jane E.~Doe\\
  Department of Electrical Engineering and Computer Sciences\\
  University of California Berkeley 
  United States\\
  \texttt{janedoe@berkeley.edu} \\
  %% examples of more authors
  %% \And
  %% Coauthor \\
  %% Affiliation \\
  %% Address \\
  %% \texttt{email} \\
  %% \AND
  %% Coauthor \\
  %% Affiliation \\
  %% Address \\
  %% \texttt{email} \\
  %% \And
  %% Coauthor \\
  %% Affiliation \\
  %% Address \\
  %% \texttt{email} \\
  %% \And
  %% Coauthor \\
  %% Affiliation \\
  %% Address \\
  %% \texttt{email} \\
}


\begin{document}
\maketitle

%===============================================================================

\begin{abstract}
    The purpose of this document is to provide both the basic paper template and submission guidelines. Abstracts should be a single paragraph, between 4--6 sentences long, ideally. Gross violations will trigger corrections at the camera-ready phase.
\end{abstract}

% Two or three meaningful keywords should be added here
\keywords{CoRL, Robots, Learning} 

%===============================================================================

\section{Introduction}
	
    Submission to CoRL 2026 will be entirely electronic, via a web site (not email). Information about the submission process and \LaTeX{} templates are available on the conference web site at \url{https://corl.org/}. For camera ready submission, use the \texttt{final} option for the \texttt{\textbackslash usepackage} command. 

%===============================================================================

\section{Citations}
\label{sec:citations}

	Citations can be made using either \textbackslash citep\{\} or \textbackslash citet\{\}, depending from the appropriateness. To avoid the citation moving to the next line, it is often a good practice to replace the space before with a tilde (\~{}) character.
	Example 1: ``CoRL is the best conference ever~\citep{Gauss1857}.''
	Example 2: ``\citet{Lagrange1788} proved, both theoretically and numerically, that CoRL is the best conference ever.''
	
%===============================================================================

\section{Experiments}
\label{sec:result}

We organize our empirical study around three questions that progressively
stress the proposed agentic VLA framework. (i) \emph{Closed-world
composition}: can an LLM-in-the-loop agent, restricted to a fixed primitive
vocabulary, solve standard manipulation benchmarks at high reliability and
package its solutions as re-usable skills? (ii) \emph{Robustness to language
and scene perturbations}: when the same task is presented with paraphrased
instructions, distractor objects, or altered visual conditions, does the
agentic decomposition recover the semantic-understanding deficit observed in
monolithic VLA baselines? (iii) \emph{Cross-environment transfer}: does the
framework generalize beyond LIBERO to other embodiments and controllers?
Section~\ref{sec:exp:setup} describes the common experimental setup and
Sections~\ref{sec:exp:libero}--\ref{sec:exp:transfer} address the three
questions in turn.

%-------------------------------------------------------------------------
\subsection{Experimental Setup}
\label{sec:exp:setup}

	\paragraph{Frozen primitive vocabulary.} To prevent the agent from
	trivializing a task by extending the action API on demand, we fix the
	primitive set throughout all experiments. The vocabulary consists of
	five core primitives---\textsc{move\_to}, \textsc{pi0\_pick},
	\textsc{release}, \textsc{set\_gripper}, and \textsc{reset}---together
	with four extended primitives that cover articulated-joint actuation
	and joint-space motion:
	\textsc{rotate\_wrist}, \textsc{rotate\_pitch},
	\textsc{js\_move\_to}, and \textsc{articulate\_to}.
	The LLM agent is permitted to compose these primitives via JSON
	commands, but is \emph{not} permitted to define new primitives during
	deployment. This constraint ensures that any reported success
	generalizes through composition rather than through ad-hoc extension
	of the action interface.

	\paragraph{Backbone VLA.} For all experiments that involve a
	visuomotor grasping primitive, we use a Pi0.5 checkpoint trained on
	the full LIBERO-130 task pool (\textsc{pi05\_libero130\_fullshot})
	as the grasping VLA. The agent invokes the VLA only through
	\textsc{pi0\_pick}, which is hard-cut after the target object lifts
	by a configurable threshold; all other motion is scripted by the
	LLM.

	\paragraph{Compliance audits.} Every successful rollout is recorded
	as a self-contained audit file containing the full command sequence,
	state and image logs, and pick/place diagnostics. We use these
	records to verify that
	(a)~the VLA call in each rollout exits before the goal predicate
	fires, and
	(b)~the goal predicate is fired by a subsequent scripted release or
	articulation action.
	A run that violates either condition is reclassified and reported
	separately rather than counted as a strict success.

%-------------------------------------------------------------------------
\subsection{LIBERO: Closed-World Skill Composition}
\label{sec:exp:libero}

	\paragraph{Benchmark.} We evaluate on the four standard LIBERO
	suites~\citep{Gauss1857}---\textsc{Spatial}, \textsc{Object},
	\textsc{Goal}, and \textsc{LIBERO-10}---each containing ten
	manipulation tasks instantiated from a shared kitchen/living-room
	scene library.

	\paragraph{Single-seed solution coverage.} For each of the 40
	(suite, task) pairs we have the agent interactively discover a
	command sequence on seed~0 of the official task ordering.
	Table~\ref{tab:libero_main} summarizes the results.
	Across all four suites, the agentic framework solves
	$40/40$ tasks under the audit criteria above. Two of the
	LIBERO-10 successes use a permitted ``doubled'' regime in which
	the VLA additionally executes a non-grasp skill (knob turn, drawer
	close) that cannot be reproduced in OSC space; the remaining 38
	rollouts are strict, meaning the VLA only performs the grasp and
	every other primitive call comes from the LLM.

	\paragraph{Skill caching and multi-seed evaluation.} A central
	claim of our framework is that exploration cost is amortized: once
	the agent has produced a working command sequence on one seed,
	the sequence should generalize to the remaining seeds of the same
	task with only the per-seed object coordinates substituted in.
	We compile each interactive solution into a Python callable
	$\pi_\text{task}: s_0 \mapsto \mathcal{C}$ that consumes the
	initial privileged state $s_0$ (object positions, articulation
	angles) and emits a list of JSON commands $\mathcal{C}$.
	The agent is no longer permitted to reason at deployment time;
	only the cached callable is invoked.
	Table~\ref{tab:libero_seeds} reports the across-seed success rate
	of these cached skills on all $50$ official initial states per
	task. The skills retain a success rate of $[\textit{XX.X}]\%$ on
	the held-out seeds, validating the hypothesis that LLM-discovered
	primitive compositions, conditioned on the runtime privileged
	state, are reproducible without further model-in-the-loop
	reasoning.

%-------------------------------------------------------------------------
\subsection{LIBERO-Pro and LIBERO-Plus: Robustness to Perturbations}
\label{sec:exp:robustness}

	A frequent failure mode of fullshot-trained VLAs is that they
	overfit to the surface statistics of the training distribution:
	even when the goal predicate is unchanged, paraphrased
	instructions, distractor objects, or altered visual textures can
	collapse the policy's success rate. The
	LIBERO-Pro~\citep{Lagrange1788} and LIBERO-Plus~\citep{Lagrange1788}
	benchmarks operationalize this by exposing each base task to
	systematic perturbations along several axes:

	\begin{itemize}
	\item \textbf{Language perturbations.} Instructions are
	rewritten with verbose paraphrases, irrelevant context, or
	stylistic variations while preserving the underlying goal.
	\item \textbf{Object/scene perturbations.} Distractor objects
	are added to the scene; existing objects are repositioned outside
	the training-state distribution; and visual conditions
	(lighting, table texture, background) are altered.
	\end{itemize}

	\paragraph{Why the agentic decomposition is expected to help.}
	The monolithic Pi0.5 backbone has been observed to behave in a
	prompt-blind manner on LIBERO-Spatial: substituting the language
	instruction with a semantically incorrect string produces
	near-identical action distributions, indicating that the
	policy's text channel is largely vestigial.
	The agentic decomposition replaces the language channel
	functionally: the LLM parses the perturbed instruction, extracts
	the canonical target object and goal predicate, and then issues
	a \textsc{pi0\_pick} call whose prompt and \texttt{track\_obj}
	argument are derived from the canonical task rather than from
	the perturbed surface form. Visual perturbations remain
	challenging because they enter the VLA through the image
	channel; we therefore expect a clean win on language and
	distractor axes and a partial win on lighting/texture.

	\paragraph{Evaluation protocol.} For each base task we evaluate
	a fixed budget of perturbed instances per perturbation axis, with
	the same cached agentic skill that succeeded on the unperturbed
	task. We report success rate against the same Pi0.5 backbone
	executed end-to-end as the baseline. Table~\ref{tab:libero_pro}
	and Table~\ref{tab:libero_plus} provide the per-axis
	breakdowns; Figure~\ref{fig:robustness_curve} aggregates the
	results into success-vs-perturbation-strength curves.

%-------------------------------------------------------------------------
\subsection{Beyond LIBERO: Cross-Environment Generalization}
\label{sec:exp:transfer}

	The final experimental phase tests whether the agentic framework
	transfers to environments that share neither LIBERO's action
	space nor its primitive set. We evaluate on two contrasting
	benchmarks:

	\begin{itemize}
	\item \textbf{CALVIN}~\citep{Gauss1857}, a single-arm
	long-horizon language-conditioned manipulation benchmark with a
	7-DoF controller that differs from LIBERO's OSC interface and a
	separate set of articulated kitchen fixtures.
	\item \textbf{RoboTwin}~\citep{Lagrange1788}, a bimanual
	manipulation benchmark whose primitive vocabulary must include
	bimanual coordination (object hand-off, two-handed grasps)
	absent from the LIBERO primitive set.
	\end{itemize}

	\paragraph{Primitive-set transfer cost.} Because the LIBERO
	primitive vocabulary is frozen, we must re-implement the
	corresponding primitives against each new controller before
	the agent can act. We report the lines-of-code overhead of this
	re-implementation in Table~\ref{tab:primitive_transfer} as a
	proxy for engineering cost, distinguishing primitives that
	transferred unmodified (e.g.\@ release, set\_gripper) from
	those that required controller-specific reimplementation
	(\textsc{move\_to}, \textsc{articulate\_to}).

	\paragraph{Task-level generalization.} Once the new primitive
	set is in place, we measure whether the agent's compositional
	pattern (frozen vocabulary, LLM-in-the-loop, VLA-only-for-grasp)
	survives the embodiment change. Table~\ref{tab:cross_env}
	reports success rates on a matched subset of CALVIN and
	RoboTwin tasks against the corresponding monolithic baselines
	on each benchmark.

	\paragraph{Limitations.} For environments lacking a publicly
	released VLA backbone, the \textsc{pi0\_pick} primitive is
	replaced by an analytic grasp primitive
	(\textsc{ik\_grasp}); this isolates the contribution of the
	LLM-side scheduling from the contribution of the VLA grasping
	model and is reported as a separate row in
	Table~\ref{tab:cross_env}.

%===============================================================================

\section{Conclusion}
\label{sec:conclusion}

	Nam dui ligula, fringilla a, euismod sodales, sollicitudin vel, wisi. Morbi auctor lorem non justo.
	Nam lacus libero, pretium at, lobortis vitae, ultricies et, tellus. Donec aliquet, tortor sed accumsan
	bibendum, erat ligula aliquet magna, vitae ornare odio metus a mi. Morbi ac orci et nisl hendrerit
	mollis. Suspendisse ut massa. Cras nec ante. Pellentesque a nulla. Cum sociis natoque penatibus
	et magnis dis parturient montes, nascetur ridiculus mus. Aliquam tincidunt urna. Nulla ullamcorper
	vestibulum turpis. Pellentesque cursus luctus mauris.
	Nulla malesuada porttitor diam. Donec felis erat, congue non, volutpat at, tincidunt tristique, libero.
	Vivamus viverra fermentum felis. Donec nonummy pellentesque ante. Phasellus adipiscing semper
	elit. Proin fermentum massa ac quam. Sed diam turpis, molestie vitae, placerat a, molestie nec, leo.
	Maecenas lacinia. Nam ipsum ligula, eleifend at, accumsan nec, suscipit a, ipsum. Morbi blandit
	ligula feugiat magna. Nunc eleifend consequat lorem. Sed lacinia nulla vitae enim. Pellentesque tin-
	cidunt purus vel magna. Integer non enim. Praesent euismod nunc eu purus. Donec bibendum quam
	in tellus. Nullam cursus pulvinar lectus. Donec et mi. Nam vulputate metus eu enim. Vestibulum
	pellentesque felis eu massa.
	Quisque ullamcorper placerat ipsum. Cras nibh. Morbi vel justo vitae lacus tincidunt ultrices. Lorem
	ipsum dolor sit amet, consectetuer adipiscing elit. In hac habitasse platea dictumst. Integer tempus
	convallis augue. Etiam facilisis. Nunc elementum fermentum wisi. Aenean placerat. Ut imperdiet,
	enim sed gravida sollicitudin, felis odio placerat quam, ac pulvinar elit purus eget enim. Nunc vitae
	tortor. Proin tempus nibh sit amet nisl. Vivamus quis tortor vitae risus porta vehicula.

%===============================================================================

\clearpage
% The acknowledgments are automatically included only in the final and preprint versions of the paper.
\acknowledgments{If a paper is accepted, the final camera-ready version will (and probably should) include acknowledgments. All acknowledgments go at the end of the paper, including thanks to reviewers who gave useful comments, to colleagues who contributed to the ideas, and to funding agencies and corporate sponsors that provided financial support.}

%===============================================================================

% no \bibliographystyle is required, since the corl style is automatically used.
\bibliography{example}  % .bib

\end{document}
